\section{问题二：碳化硅外延层的联合反演及线形核验}
\label{sec:sic-results}

\subsection{确定可使用的色散窗口}
附件1与附件2测量同一晶圆，分别对应 $10^\circ$ 和 $15^\circ$ 入射。
两条全谱在声子区均出现高反射平台，而高波数区呈现衰减条纹。
本文据式\eqref{eq:wang-sic}的明确适用范围选取
$2000$--$3300\,\mathrm{cm}^{-1}$ 为主窗口，对应波长约为
$5.000$--$3.030\,\mu\mathrm m$，从而避开强声子响应，并避免将
折射率公式外推到其 $5\,\mu\mathrm m$ 上限之外\cite{Wang2013,Polyanskiy2024}。
在两个角度统一剔除 $2300$--$2400\,\mathrm{cm}^{-1}$ 的局部异常段，
余下每角 $2488$ 点。该异常与大气吸收区相近，但没有参考气体谱，
故只将其作为局部污染处理，不宣称已经确定具体来源。

式\eqref{eq:wang-sic}结合 $\lambda=10^4/\sigma$ 给出的折射率
在2000、2500和3300$\,\mathrm{cm}^{-1}$处分别为
2.44397、2.49333和2.52811。相应法向群折射率
$n_g=n+\sigma\,\mathrm{d}n/\mathrm{d}\sigma$ 为
2.72756、2.66231和2.61857，说明本窗口虽远离反射平台，
仍不能把相位折射率与控制条纹周期的群折射率混用。
样品晶型和偏振未知，后文几何厚度均以这一普通光色散为条件。
\begin{table}[htbp]\centering
\caption{碳化硅主窗口内的文献折射率及群折射率}\label{tab:sic-index}
\begin{tabular}{lrr}
\toprule
波数/$\mathrm{cm}^{-1}$ & $n$ & $n+\sigma\,\mathrm{d}n/\mathrm{d}\sigma$\\\midrule
2000 & 2.443966 & 2.727559\\
2500 & 2.493330 & 2.662306\\
3300 & 2.528109 & 2.618570\\
\bottomrule
\end{tabular}\end{table}

表\ref{tab:sic-index}保留折射率计算的中间值，便于检查波长单位与色散导数。

\subsection{FFT 与峰级次给出独立初值}
先在未挖去局部区间的完整主窗口内扣除三次趋势，乘 Hann 窗后计算 FFT，
在 $0.002$--$0.008\,\mathrm{cm}$ 的频率区间寻找主峰。
两角均得到约 $0.003910\,\mathrm{cm}$，对应周期
$255.776\,\mathrm{cm}^{-1}$。若暂取 $n=2.55$ 并补入各自的入射角，
基准厚度分别为 $7.684$ 和 $7.706\um$。
这里频率单位为波数单位的倒数；$2^{18}$ 点零填充只细化峰值网格，
原始窗口决定的频率尺度仍约为 $7.70\times10^{-4}\,\mathrm{cm}$。
因此，这两个数用于定位搜索范围，不作为高精度结果。

另一条初值路径使用 $31$ 点三次 Savitzky--Golay 平滑定位局部极大值\cite{SavitzkyGolay1964}，
然后恢复相对干涉级次。附件1保留的峰为
$2078.887,2584.627,2830.025,3094.225\,\mathrm{cm}^{-1}$，
附件2为 $2101.547,2608.733,2870.523,3122.670\,\mathrm{cm}^{-1}$。
由于异常区间覆盖一个条纹峰，两组对应的级次均为 $(0,2,3,4)$，
不能将跨越缺口的两个峰误当成相邻级次。
将这些级次与色散坐标 $X_\theta(\sigma)$ 回归，得到
$7.442$ 和 $7.391\um$。其与后续结果处在同一厚度分支，
但缓变背景造成的峰位偏移仍需由全谱拟合处理。
\begin{table}[htbp]\centering
\caption{碳化硅峰位与恢复的相对条纹级次}\label{tab:sic-peaks}
\begin{tabular}{lrr}
\toprule
相对级次 & 附件1峰位/$\mathrm{cm}^{-1}$ & 附件2峰位/$\mathrm{cm}^{-1}$\\\midrule
0 & 2078.887 & 2101.547\\
2 & 2584.627 & 2608.733\\
3 & 2830.025 & 2870.523\\
4 & 3094.225 & 3122.670\\
\bottomrule
\end{tabular}\end{table}

表\ref{tab:sic-peaks}中的级次跳跃与被屏蔽区间相对应，计算时保留该跳跃。

\paperfigure{figures/sic_fft.png}{碳化硅条纹的 FFT 初值检验；零填充频率网格不增加独立光谱分辨率}{0.80}{fig:sic-fft}

\subsection{变量投影下的双角共享厚度拟合}
令 $t=2(\sigma-2000)/1300-1$，用三次 Chebyshev 多项式描述背景，
二次多项式描述条纹包络。每个角度的模型写成
\begin{align}
 \widehat R_j(\sigma)
 &=\sum_{k=0}^{3}b_{jk}T_k(t)
 +\left(\sum_{k=0}^{2}a_{jk}t^k\right)\cos\Phi_j(\sigma),\\
 \Phi_j(\sigma)&=4\pi\times10^{-4}d\,
 \sigma\sqrt{n^2(\sigma)-\sin^2\theta_j}+\phi_j,
 \qquad j=1,2.
 \label{eq:sic-regression}
\end{align}
其中 $R_j$ 直接以反射率百分数计算，残差单位为百分点。
两个角度共享 $d$，分别保留常数相移 $\phi_j$，容纳界面反射相位及
有限窗口内的有效偏移。背景和包络只提供低阶缓变项，不自由逐点追踪条纹。

给定非线性参数 $\boldsymbol\eta=(d,\phi_1,\phi_2)$ 后，
式\eqref{eq:sic-regression}对每角的七个系数均为线性模型。
构造设计矩阵 $M_j(\boldsymbol\eta)$，先解
\begin{equation}
 \widehat{\boldsymbol c}_j(\boldsymbol\eta)
 =\operatorname*{argmin}_{\boldsymbol c_j}
 \|M_j(\boldsymbol\eta)\boldsymbol c_j-\boldsymbol R_j\|_2^2,
 \quad
 \widehat{\boldsymbol\eta}
 =\operatorname*{argmin}_{\boldsymbol\eta}
 \sum_{j=1}^{2}\|M_j\widehat{\boldsymbol c}_j-\boldsymbol R_j\|_2^2.
 \label{eq:sic-variable-projection}
\end{equation}
这一变量分离思路\cite{GolubPereyra1973}将外层非线性搜索限制为三个参数，避免将全部背景系数一并交给
非线性优化器。搜索约束为 $5\leq d\leq10\um$；在每三点取一点的谱上，
以 $6.6$--$8.0\um$、间隔 $0.2\um$ 的厚度网格和不同相位启动搜索，
随后用全部保留原始点精化。平滑谱只参与峰值初值，主回归不使用平滑后的反射率。
每次目标函数调用均重新解线性最小二乘，而不是先固定一条扣背景曲线再拟合相位，
使背景估计误差可以在有限自由度内与相位共同调整。两角使用相同的波数区间及等权残差；
题目没有提供逐点测量方差，故不人为构造精密的逆方差权重。
相移具有 $2\pi$ 周期，表述时将其归一化到同一主值区间；它只确定本模型内的
条纹起点，不能单凭其数值推出界面的微观组成或绝对干涉级次。

联合估计为
\begin{equation}
 \boxed{\widehat d_{\mathrm{SiC}}=7.384314\um\approx7.38\um},
 \qquad \mathrm{RMSE}=0.024400\text{ 个百分点}.
 \label{eq:sic-final}
\end{equation}
对应 $R^2=0.997640$，两角有效相移为 $3.10774$ 和 $2.65491\,\mathrm{rad}$。
\figref{fig:sic-fit}表明主条纹可由共同厚度描述；
\tabref{tab:sic-angles}同时给出分角结果，防止联合平均掩盖不一致性。
两角差为 $0.047031\um$，占联合厚度约 $0.637\%$，
超过条件残差重采样所体现的局部波动尺度，提示仍存在光学常数、
反射相位或仪器响应的系统偏差。不能仅凭较高 $R^2$ 宣称已获得同等级的绝对精度。

\begin{table}[htbp]
\centering
\caption{碳化硅分角估计与共享厚度估计}
\label{tab:sic-angles}
\begin{tabular}{lrrr}
\toprule
拟合数据 & 有效点数 & $d\,(\mu\mathrm m)$ & RMSE（百分点）\\
\midrule
附件1，$10^\circ$ & 2488 & 7.407811 & 0.015471\\
附件2，$15^\circ$ & 2488 & 7.360780 & 0.027416\\
两角共享厚度 & 4976 & 7.384314 & 0.024400\\
\bottomrule
\end{tabular}
\end{table}

\paperfigure{figures/sic_fit.png}{碳化硅双角联合拟合；阴影区为未参与回归的局部异常波段}{0.78}{fig:sic-fit}
\begin{table}[htbp]
\centering
\caption{碳化硅同一有效窗口内的折射率及线形模型比较}
\label{tab:sic-models}
\begin{tabular}{lrrl}
\toprule
模型 & $d\,(\mu\mathrm m)$ & RMSE（百分点） & 解释\\
\midrule
常数 $n=2.55$，两束 & 7.712680 & 0.025149 & 常折射率基准\\
Lorentz 透明近似，两束 & 7.336646 & 0.025007 & 替代色散情景\\
Wang 普通光，两束 & 7.384314 & 0.024400 & 主参照模型\\
Wang 普通光，有效 Airy & 7.384155 & 0.024330 & 嵌套线形检验\\
\bottomrule
\end{tabular}
\par\smallskip
\begin{minipage}{0.92\textwidth}\small
所有行均使用 $2000$--$3300\,\mathrm{cm}^{-1}$，剔除
$2300$--$2400\,\mathrm{cm}^{-1}$，并联合拟合两个角度。
Lorentz 情景取 $\epsilon_\infty=6.7$、$\sigma_{\rm TO}=797\,\mathrm{cm}^{-1}$、
$\sigma_{\rm LO}=970\,\mathrm{cm}^{-1}$，只作模型敏感性比较。
额外小数用于复核计算，不代表同等量级的绝对测量精度。
\end{minipage}
\end{table}

表中常折射率回归给出 $7.712680\um$，较主模型高 $4.447\%$。
这一差异来自在有限波段内采用不同相位梯度，并不是单位换算导致的误差。
采用透明区 Lorentz 近似
$n^2=6.7(970^2-\sigma^2)/(797^2-\sigma^2)$ 时，
厚度为 $7.336646\um$，较主模型低 $0.646\%$。
这两种对照的训练残差仍然较小，表明良好拟合本身不足以唯一确定材料色散与几何厚度。
选择 Wang 模型作为主参照的依据是可追溯的色散关系及其适用域，
不是事后挑选能产生某一预期厚度的公式。

峰级次、FFT 和非线性回归分工明确：前两者验证厚度的数量级及条纹分支，
回归利用全部有效谱点估计相位梯度。峰位受背景导数影响，而 FFT 还受有限窗口和
变周期条纹展宽影响，因此不将三个估计机械平均以制造额外精度。
主结果的复核应直接比较同一色散假设下的预测谱、残差及分角厚度，
而不是只检查最终数值是否接近。

\subsection{问题三补充：碳化硅是否需要高阶往返修正}
为在同一数据和背景自由度下检验非正弦线形，将式\eqref{eq:sic-regression}
中的余弦替换为有效 Airy 谐波核
\begin{equation}
 H(\Phi;\xi)=\frac{\cos\Phi-\xi}{1-2\xi\cos\Phi+\xi^2}
 =\sum_{m=1}^{\infty}\xi^{m-1}\cos(m\Phi),\qquad |\xi|<1.
 \label{eq:sic-airy-kernel}
\end{equation}
该式与常系数透明 Airy 强度经过背景及振幅重参数化后具有相同线形结构，
且 $H(\Phi;0)=\cos\Phi$，因而严格嵌套两束模型。
为避免与式\eqref{eq:snell-q}的纵向传播因子混淆，正文将代码中的有效
\texttt{q} 参数记为 $\xi$。每角独立估计 $\xi_j\in[-0.45,0.45]$，
并固定中心处基本谐波振幅为正来消除符号二义性。
这些参数描述相对高次谐波，不等于从数据独立测得的界面反射率。

拟合得到 $\xi_{10}=-0.003265$、$\xi_{15}=0.005908$，厚度为
$7.384155\um$，比两束结果仅小 $0.000158\um$，相对变化为
$-0.00214\%$。训练 RMSE 从 $0.024400$ 降至 $0.024330$ 个百分点，
按原始点数计算的 BIC 下降 $11.55$，但密集谱点的相关性使该指标不能单独裁决模型。

将主窗口按 $50\,\mathrm{cm}^{-1}$ 连续小块编号，再以块编号模 $5$ 分配验证折，
每折只在其他四组上拟合。两束与有效 Airy 的合并留出 RMSE 分别为
$0.032025$ 和 $0.033067$ 个百分点，增加高阶参数反而恶化约 $3.25\%$。
因此，本窗口没有提供足够稳定的证据支持替换两束主模型，保留式\eqref{eq:sic-final}
作为条件厚度估计。高阶修正远小于分角差和折射率情景变化；
这一结论仅适用于所选透明窗口，并不等价于证明全谱中没有多次反射。

\paperfigure{figures/sic_airy_cv.png}{碳化硅两束与有效 Airy 模型的连续谱块验证；训练误差的轻微改善未转化为留出预测改善}{0.78}{fig:sic-airy-cv}
